\documentclass[12pt,a4paper,twoside]{article}

% are used to comment out text

%------------------------------------------------------------------------------------------------------------
%EDIT THESE FIELDS WITH YOUR DETAILS
%------------------------------------------------------------------------------------------------------------
\newcommand{\studentname}{}
%\newcommand{\studentnumber}{}
\title{Experimento de Validação do Dispositivo de Calibração Eletrostática (DCE) \\ Roteiro}
\newcommand{\tutor}{}
%\newcommand{\labgroup}{}
\newcommand{\labdate}{July 2026}
\newcommand{\reportdate}{XX Julho de 2026}
%------------------------------------------------------------------------------------------------------------

%------------------------------------------------------------------------------------------------------------
%PACKAGES. PACKAGES ARE USED TO EXTEND THE FUNCTIONALITY OF LATEX. YOU SHOULDN'T
%NEED ANY MORE, SPEND YOUR TIME ON THE CONTENT!
%------------------------------------------------------------------------------------------------------------

\usepackage[a4paper,left=25mm,right=25mm,bottom=25mm,top=25mm]{geometry}
\usepackage{fancyhdr}
\usepackage{subfigure}
\usepackage{graphicx}
\usepackage{cleveref}
\usepackage{booktabs}
\usepackage[british]{babel}
\usepackage[square,comma,numbers,sort&compress]{natbib}
\usepackage{csvsimple}
\usepackage{graphicx}
\usepackage{pgfplotstable}
\usepackage{gensymb}
\usepackage{array}
\usepackage{multirow}
\usepackage{url}
\pgfplotsset{compat=1.9}% supress warning
\begin{document}

%------------------------------------------------------------------------------------------------------------

%------------------------------------------------------------------------------------------------------------
%DO NOT EDIT THESE FIELDS
%------------------------------------------------------------------------------------------------------------
\setlength{\parindent}{0em}
\setlength{\parskip}{0.5em}
%\author{\studentname \\ Orientação: \tutor.}
\date{Data do Relatório: \reportdate.}

\pagestyle{fancy}
\fancyhead{}
%\fancyhead[C]{\studentname~\studentnumber}
\fancyfoot{}
\fancyfoot[C]{\thepage}
\renewcommand{\headrulewidth}{0pt}
\renewcommand{\footrulewidth}{0pt}

\renewcommand\bibname{Referências}

%------------------------------------------------------------------------------------------------------------

\maketitle

%------------------------------------------------------------------------------------------------------------
% REPORT BEGINS HERE
%------------------------------------------------------------------------------------------------------------
%%%%%%%%%%%%%%%%
\section{Protocolos de Segurança e EPIs}

Antes de iniciar qualquer etapa prática, os procedimentos de segurança contra Alta Tensão devem ser estritamente seguidos.

\subsection*{Equipamentos de Proteção Individual (EPIs):}
\begin{itemize}
    \item \textbf{Luvas Isolantes (Classe 0):} Específicas para trabalhos elétricos, certificadas para suportar até 1000\,V AC ou 1500\,V DC. Devem ser usadas com luvas de vaqueta por cima para proteção mecânica.
    \item \textbf{Calçado de Segurança com Isolamento Elétrico:} Calçado fechado sem componentes metálicos.
    \item \textbf{Óculos de Segurança:} Essencial caso ocorra a ruptura dielétrica do ar (centelha) entre as placas do DCE, o que pode ejetar pequenas partículas de alumínio.
    \item \textbf{Jaleco/Vestimenta 100\% Algodão:} Materiais sintéticos derretem na pele em caso de arco elétrico.
\end{itemize}

\subsection*{Equipamentos de Proteção Coletiva (EPCs) e Cuidados Especiais:}
\begin{itemize}
    \item \textbf{Tapete Isolante de Borracha:} Deve estar posicionado no chão, onde o operador da fonte HV ficará em pé.
    \item \textbf{Bastão de Aterramento (Bastão de Descarga):} Essencial para o pós-teste. O DCE atua como um capacitor e armazena carga. Mesmo com a fonte desligada, as placas podem reter carga suficiente para causar um choque severo.
    \item \textbf{Sinalização e Isolamento:} A área ao redor da balança deve estar demarcada com faixas de ``Perigo: Alta Tensão''.
    \item \textbf{Prevenção de Ruptura Dielétrica:} Centelhas aparecem na extremidade de condutores quando o campo elétrico provoca a ruptura dielétrica do ar. A 1000\,V em pressão atmosférica, se as placas estiverem muito coladas (ex: menor que 0,5\,mm), uma centelha pode ocorrer, podendo fechar um curto-circuito na sua fonte. Garanta que a distância (\textit{gap}) seja segura.
\end{itemize}

\section{Preparação do Sistema (Com a fonte totalmente desligada)}

\begin{enumerate}
    \item \textbf{Montagem Mecânica:} Fixe um disco de alumínio no braço móvel do pêndulo e o outro disco no suporte fixo da base.
    \item \textbf{Alinhamento:} Utilize um paquímetro ou micrômetro para ajustar a distância $D_E$ (entre-ferro) inicial entre os discos de alumínio. Garanta que as placas estejam perfeitamente paralelas.
    \item \textbf{Cabeamento:} Conecte os cabos de alta tensão da fonte de tensão Magna-Power aos eletrodos do DCE. Garanta que os cabos estejam fixos e não interfiram excessivamente na rigidez mecânica (constante elástica) do pêndulo.
    \item \textbf{Aterramento:} Verifique se o chassi da balança e a carcaça da fonte de tensão estão devidamente aterrados na malha do laboratório.
\end{enumerate}

\section{Execução do Experimento (Teste Estático em Atmosfera)}

Este procedimento descreve a calibração primária do DCE, onde a constante $K_{dce}$ será validada na pressão atmosférica, antes de ir para o vácuo.

\begin{enumerate}
    \item \textbf{Afastamento:} Todos os envolvidos devem se afastar da balança e ficar atrás de uma barreira ou a uma distância segura, sobre o tapete isolante.
    \item \textbf{Ligação do Sistema:} Ligue a alimentação geral do laboratório e ligue a fonte Magna-Power, garantindo que o controle de tensão esteja zerado (0\,V) e o limite de corrente (corrente de corte) configurado para um valor mínimo de segurança (ex: 10\,mA, pois o DCE não consome corrente em estado estacionário por ser um capacitor; se puxar corrente, há um curto ou fuga).
    \item \textbf{Degraus de Tensão:} Aumente a tensão para 100\,V. \\
    \textit{Atenção:} Ao aplicar a tensão contínua, o pêndulo será atraído e passará a oscilar em torno de uma nova coordenada de equilíbrio.
    \item \textbf{Regime Estacionário:} É estritamente necessário aguardar a completa atenuação do regime transiente (oscilações) até que o sistema atinja o regime estacionário para efetuar a leitura do deslocamento. Aguardar 3 minutos para garantir que as oscilações tenham cessado.
    \item \textbf{Registro de Dados:} Salve o gráfico de deslocamento do pêndulo (eixo Y) versus tempo (eixo X) usando o software de aquisição. Anote o valor do deslocamento final (posição de equilíbrio) e a tensão aplicada obtido pelo script para cada caso.
    \item \textbf{Repetitibilidade}: Repita o passo de 100\,V mais 4 vezes para garantir a repetibilidade dos dados.
    \item \textbf{Progressão:} Repita os passos de aumento de tensão em incrementos de 100\,V (100\,V, 200\,V, 300\,V, 400\,V, 500\,V, 600\,V, 700\,V, 800\,V e 1000\,V). Espere o regime estacionário e anote os dados em cada degrau.
\end{enumerate}

\section{Desenergização e Coleta de Dados}

\begin{enumerate}
    \item \textbf{Redução de Tensão:} No painel da fonte Magna-Power, reduza gradativamente a tensão de 1000\,V de volta para 0\,V.
    \item \textbf{Desligamento Geral:} Desligue o botão de \textit{Output} da fonte e, em seguida, a chave geral da fonte.
    \item \textbf{Procedimento de Descarga (Crítico):} 
    \begin{itemize}
        \item Coloque as Luvas Isolantes Classe 0.
        \item Pegue o bastão de aterramento.
        \item Encoste a ponta do bastão simultaneamente nos terminais de ambas as placas do DCE (ou faça o aterramento de uma placa contra a outra) para curtocircuitar o capacitor e garantir que toda a energia potencial armazenada no campo elétrico seja dissipada.
        \item Mantenha o bastão no local por cerca de 5 segundos.
    \end{itemize}
    \item \textbf{Análise de Dados:}
    \begin{itemize}
        \item Organize os dados coletados em um gráfico. No eixo X, coloque o quadrado da tensão aplicada ($V^2$) e no eixo Y, o deslocamento medido ou a força calculada.
        \item A relação deverá ser estritamente linear, comprovando a equação teórica da força do capacitor de placas paralelas ($F_c \propto V^2/d^2$). A inclinação desta reta fornecerá a constante experimental do seu DCE.
    \end{itemize}
\end{enumerate}

%%%%%%%%%%%%%%%%
\section{Resultados}


%%%%%%%%%%%%%%%%
\section{Conclusões}


\newpage
\bibliographystyle{unsrtnat}
%\bibliography{refs}

\end{document}
